\section{مدل های \lr{Score-Based}}
\subsection{سوال اول}
از جمله مشکلاتی که الگوریتم \lr{Langevin Dynamics} ممکن است به آنها بربخورد عبارتند از :
\begin{enumerate}
    \item \lr{Slow Mixing} : ایده اصلی این الگوریتم حرکت به سمت نواحی با احتمال بالاتر با استفاده از اطلاعات تابع \lr{score} هست، ولی در فضاهای با ابعاد بالا این پروسه میتواند به شدت آرام باشد و به دلیل استفاده از نویز و استوکسنیک بودن ممکن است الگوریتم در مینیموم های محلی گیر کند.
    \item \lr{Noise Schedule} : الگوریتم \lr{langevin Dynamics} به شدت به نویز حساس هست و باید تنظیم مقدار نویز با دقت انجام شود.
\end{enumerate}
روش هایی که برای بهبود این الگوریتم میتوانیم استفاده کنیم عبارت اند از :
\begin{enumerate}
    \item کاهش واریانس :  با استفاده از تکنیک‌هایی مانند نمونه‌گیری مهم و یا روش‌های بازتوزیع‌سازی مبتنی بر گرادیان (\lr{score}). این تکنیک‌ها باعث کاهش واریانس فرآیند نمونه‌برداری و امکان استفاده از گام‌های بزرگتر و همگرایی سریع‌تر می‌شوند.
    \item \lr{Adaptive Step Size} : الگوریتم‌هایی که به صورت خودکار اندازه گام را بر اساس میزان انحنای محلی تابع گرادیان تنظیم می‌کنند. در مناطق با انحنای بالا، گام کاهش می‌یابد و در مناطق صاف، گام افزایش می‌یابد.
    \item \lr{SGLD} :  این روش یک نسخه محبوب از لَنگَوین داینامیک است که تنها از زیرمجموعه‌ای از داده‌ها (یک دسته کوچک) برای محاسبه گرادیان در هر مرحله استفاده می‌کند. این کار باعث کاهش قابل توجه هزینه محاسباتی و امکان نمونه‌برداری از مجموعه‌های داده بسیار بزرگ می‌شود.
    \item \lr{Noisy Neural Langevin Dynamics} :  این رویکرد از قدرت شبکه‌های عصبی برای یادگیری یک برنامه نویز آموزنده‌تر استفاده می‌کند، که منجر به بهبود کارایی نمونه‌برداری می‌شود.
\end{enumerate}

\subsection{سوال دوم}
برای ساخت یک مدل شرطی برای تولید نمونه از کلاس‌های مختلف، می‌توان از یک مدل score-based غیر شرطی از پیش آموزش دیده به همراه یک طبقه‌بند نرم \lr{(soft classifier)} استفاده کرد. هدف این است که مدل \lr{score-based} را با اطلاعات طبقه‌بندی یاد گرفته‌شده، هدایت کرده و نمونه‌هایی تولید کنیم که به احتمال زیاد به کلاس‌های مورد نظر تعلق داشته باشند.

ایده اصلی این است که از گرادیان تابع لگاریتم احتمال \lr{(score function)} مدل \lr{score-based} غیر شرطی استفاده کنیم و آن را با یک ترم هدایت‌کننده \lr{(guidance term)} بر اساس خروجی طبقه‌بند نرم ترکیب کنیم. این ترم هدایت‌کننده، نمونه‌ها را به سمت فضاهایی که احتمال تعلق به یک کلاس خاص بیشتر است، سوق می‌دهد.


فرض کنید:
\begin{enumerate}
    \item   $x_t$ نشان‌دهنده نمونه در زمان $t$ است.
    \item $\epsilon$ نرخ نویز \lr{(noise schedule)} است.
    \item  $s(x_t) = \nabla_x \log p(x_t)$ گرادیان تابع لگاریتم احتمال \lr{(score function)} مدل \lr{score-based} غیر شرطی است.
    \item  $p_c(y|x_t)$ خروجی طبقه‌بند نرم برای کلاس $y$ با توجه به نمونه $x_t$ است. (یک توزیع احتمال بر روی کلاس‌ها)
    \item $\beta$ یک ضریب وزن‌دهی \lr{(weighting factor)} است که قدرت هدایت شرطی را کنترل می‌کند.
\end{enumerate} 

به‌روزرسانی نمونه در الگوریتم \lr{Langevin dynamics} به صورت زیر تغییر می‌کند:

$$
x_{t+1} = x_t + \epsilon s(x_t) + \sqrt{2\epsilon} z_t + \beta \nabla_{x} \log p_c(y|x_t)
$$

که در آن $z_t$ یک نویز تصادفی با توزیع نرمال استاندارد است.
\begin{enumerate}
\item  عبارت $\epsilon s(x_t)$ نشان‌دهنده گام تصادفی در جهت گرادیان تابع لگاریتم احتمال است.
\item  عبارت $\sqrt{2\epsilon} z_t$ نشان‌دهنده نویز است که برای جلوگیری از همگرایی سریع و امکان اکتشاف در فضای نمونه‌ها اضافه می‌شود.
\item  عبارت $\beta \nabla_{x} \log p_c(y|x_t)$  ترم هدایت‌کننده است که نمونه را به سمت فضاهایی که احتمال تعلق به کلاس $y$ بیشتر است، سوق می‌دهد.  مقدار $\beta$ تعیین می‌کند که چقدر این هدایت باید قوی باشد.
\end{enumerate}
به این ترتیب، با افزودن این ترم به الگوریتم \lr{Langevin dynamics}، می‌توانیم از مدل \lr{score-based} غیر شرطی برای تولید نمونه‌هایی استفاده کنیم که به طور خاص به یک یا چند کلاس خاص مرتبط هستند، با هدایت آن‌ها توسط خروجی طبقه‌بند نرم.

\subsection{سوال سوم}

در ساخت مدل شرطی با استفاده از روش مطرح شده در قسمت (2)، نمی‌توان از یک طبقه‌بند از پیش آموزش‌دیده یخ‌زده \lr{frozen pretrained classifier} استفاده کرد. این امر به دلایل زیر ضروری است:
\begin{enumerate}

\item  عدم تطابق توزیع‌ها: طبقه‌بند از پیش آموزش‌دیده بر روی یک مجموعه داده و یک وظیفه خاص آموزش داده شده است. خروجی‌های آن ممکن است با توزیع‌هایی که توسط مدل \lr{score-based} تولید می‌شوند، همسو نباشد. ترم هدایت شرطی به خروجی‌های طبقه‌بند برای راهنمایی فرآیند نمونه‌برداری متکی است. یک طبقه‌بند یخ‌زده این امکان را فراهم نمی‌کند.

\item  عدم همسویی با نمونه‌های تولید شده: مدل \lr{score-based} نمونه‌ها را بر اساس توزیع احتمال کلی آن تولید می‌کند. این نمونه‌ها ممکن است به طور ذاتی با کلاس‌هایی که ما علاقه‌مند به آنها هستیم، همسو نباشند. طبقه‌بند باید این نمونه‌ها را به کلاس‌های مورد نظر نگاشت کند. یک طبقه‌بند یخ‌زده این نگاشت را یاد نگرفته است.

\item  اثرگذاری بر \lr{Score Function}: حتی اگر از یک طبقه‌بند یخ‌زده استفاده کنیم، خروجی آن می‌تواند باعث ایجاد سوگیری در تابع \lr{score} شود. تابع \lr{score} باید بازتابی از توزیع احتمال واقعی داده‌ها باشد. یک طبقه‌بند یخ‌زده، آموزش‌دیده بر روی توزیع متفاوت، می‌تواند تابع \lr{score} را فاسد کند و منجر به نتایج ضعیف شود.
\end{enumerate}


طبقه‌بند شرطی، نقش حیاتی در ترجمه نمونه‌های غیر شرطی تولید شده توسط مدل \lr{score-based} به نمونه‌هایی دارد که به طور دقیق با کلاس‌های مورد نظر مطابقت دارند. این ترجمه نیازمند سازگاری و تنظیم دقیق طبقه‌بند با ویژگی‌های نمونه‌های تولید شده است، که یک طبقه‌بند یخ‌زده فاقد این قابلیت است.

\subsection{سوال چهارم}
تابع هدفی که سعی در بهینه سازی آن داریم به شکل زیر هست ؛:
\begin{equation*}
    J(\theta) = \mathbb{E}_{P_{data}} \left\lceil \frac{1}{2} ||s_\theta (x)||^2_2 +\operatorname{tr}(\nabla_x s_\theta(x))\right\rceil 
\end{equation*}
\begin{enumerate}
    \item عبارت اول به عنوان پنالتی برای انرژی هست. این عبارت اندازه تابع \lr{score} را جریمه میکند. از آنجایی که تابع اسکور به ناحیه های با احتمال بیشتر اشاره میکند این هبارت اول از به وجود آمدن پیک های تیز جلوگیری میکند.
    \item ترم دوم که درواقع لاپلاسین $\log p_{data}$ هست خم فضا رو (\lr{curvature}) را کنترل میکند تا در مد ها یا همان دیتا پوینت ها شکل تابع معقر باشد و ماکسیموم های محلی به وجود بیاید.
\end{enumerate}



